\documentclass[10pt,a4paper]{amsart}
\usepackage[margin=19mm]{geometry}
\usepackage{amsmath,amssymb,mathtools}
\usepackage[scr=rsfs,cal=euler]{mathalfa}
\usepackage{xfrac}
\usepackage[unicode,hidelinks,bookmarks=false]{hyperref}
\newcommand{\ie}{i.e.\ }
\newcommand{\cf}{cf.\ }
\newcommand{\resp}{respectively\ }
\newcommand{\Subsection}{\subsection}
\providecommand{\nobreakdash}{\nobreak\hbox{-}}
\setlength{\emergencystretch}{2em}
\multlinegap0pt
\usepackage{bm}
\usepackage{bbm}
\usepackage{dsfont}
\usepackage{xspace}
\usepackage{cancel}
\usepackage{mathtools}
\usepackage{amsthm}
\makeatletter
\@ifundefined{theorem}{\newtheorem{theorem}{Theorem}}{}
\@ifundefined{proposition}{\newtheorem{proposition}[theorem]{Proposition}}{}
\makeatother
\newcommand{\IN}{\mathbb{N}} % Natural numbers #listed
\newcommand{\KP}{\mathcal{P}}
\newcommand{\tightlist}{\setlength{\itemsep}{0pt}\setlength{\parskip}{0pt}}
\title[Discrete Fa\`a di Bruno via M\"obius Inversion]{Discrete Fa\`a di Bruno via M\"obius Inversion}
\author{Heinrich Hartmann}
\date{Source: arXiv:2607.07742v1 (CC BY 4.0)}
\begin{document}
\maketitle

\noindent\emph{Source and licence.} Discrete Fa\`a di Bruno via M\"obius Inversion. Version arXiv:2607.07742v1 (CC BY 4.0), distributed under the Creative Commons Attribution 4.0 International licence, as recorded in the arXiv licence metadata for this exact version and shown on the abstract page https://arxiv.org/abs/2607.07742v1. Full rights notice: licenses/arxiv-2607.07742-CC-BY-4.0.txt.\par\smallskip

\noindent\textit{Editorial scope.} Counted unit: the Theorem whose source label is \texttt{thm:product-rule} in the article (source0.tex lines 2254--2294, source label \texttt{thm:product-rule}), delivered together with the article's own complete original proof of it (source0.tex lines 2296--2318, one \texttt{proof} environment of 23 lines). No mathematics is written, repaired or re-derived: statement and proof are the article's own text line for line. Required context retained uncounted: prop:taylor-duality (lines 942--971). Every definition, display and label of those lines is retained unchanged. Omitted from this package and not redistributed: 1--941, 972--2253, 2295--2295, 2319--2417. Nothing inside a retained excerpt is omitted or altered: every retained body is delivered exactly as the article's lines are written, so no omission receipt of any kind accompanies this package. Citations: the retained excerpts carry no citation command at all, so no bibliography entry of the article is transcribed and no reference is redistributed. Reference adaptation: none was needed and none was made; every reference inside the delivered unit resolves inside it, so no unresolved reference or citation is produced. Printed slips kept literal: no printed word, symbol, hypothesis or number of the counted statement or of its proof was altered. Layout and class: the article's own class is \texttt{article} with packages including inputenc, fontenc, lmodern, graphicx, geometry, setspace, enumitem, tocloft, amsmath, amssymb, mathtools, amsthm, longtable, natbib; the wrapper is the lane's pinned \texttt{amsart} editorial wrapper at 10pt on A4, which restates the theorem-like environments the retained text needs and adds the article's own macro definitions that the retained text needs (\texttt{\textbackslash IN}, \texttt{\textbackslash KP}, \texttt{\textbackslash tightlist}), with extra wrapper packages (bm, bbm, dsfont, xspace, cancel, mathtools). The delivered numbering is the unit's own. Assets: no retained span contains a figure environment, a graphics-inclusion command, a PDF-page inclusion, a table or a TikZ picture; no image file is copied into this package, the package declares no asset dependency and no image file is redistributed with it. Licence: version arXiv:2607.07742v1 of this article is distributed under the Creative Commons Attribution 4.0 International licence, as recorded in the arXiv licence metadata for this exact version and shown on the abstract page https://arxiv.org/abs/2607.07742v1. A full rights notice is delivered at licenses/arxiv-2607.07742-CC-BY-4.0.txt. Characters: every retained excerpt is pure ASCII, so the delivered engine, XeLaTeX with its default Latin Modern text font, reports no missing character.
\begin{proposition}[Taylor duality]

\label{prop:taylor-duality} Let \(X, Y\) be abelian groups,
\(g: X \to Y\), \(x \in X\), \(u_\bullet = (u_1, \dots, u_k) \in X^k\),
and \(y = g(x)\).

\begin{enumerate}
\def\labelenumi{\arabic{enumi})}
\tightlist
\item
  \emph{Boolean Taylor duality.}
\end{enumerate}

\[
  g(x + {\textstyle\sum_{i=1}^k} u_i) = y + \sum_{\emptyset \neq S \subseteq [k]} \Delta(g; x; u_S), \qquad \Delta(g; x; u_\bullet) = \sum_{S \subseteq [k]} (-1)^{k-|S|}\, g(x + {\textstyle\sum_{i \in S}} u_i).
\]

\begin{enumerate}
\def\labelenumi{\arabic{enumi})}
\setcounter{enumi}{1}
\tightlist
\item
  \emph{Binomial Taylor duality.} For \(\alpha \in \IN_0^k\):
\end{enumerate}

\[
  g(x + {\textstyle\sum_{i=1}^k} \alpha_i u_i) = y + \sum_{0 < \beta \leq \alpha} \frac{(\alpha)_\beta}{\beta!}\, \Delta(g; x; u_\bullet^\beta), \qquad \Delta(g; x; u_\bullet^\alpha) = \sum_{\beta \leq \alpha} (-1)^{|\alpha|-|\beta|}\, \frac{(\alpha)_\beta}{\beta!}\, g(x + {\textstyle\sum_{i=1}^k} \beta_i u_i).
\]

\end{proposition}

\begin{theorem}[Discrete product rule]

\label{thm:product-rule} Let \(X\) be an abelian group, \(A\) an
associative algebra, and \(f_1,\dots,f_r: X \to A\). For \(x \in X\) and
\(u_1,\dots,u_n \in X\):

\begin{enumerate}
\def\labelenumi{\arabic{enumi})}
\tightlist
\item
  \emph{Boolean product rule:}
\end{enumerate}

\[
  \Delta(f_1 \cdots f_r; x; u_1,\dots,u_n)
  =
  \sum_{\substack{J_1,\dots,J_r \subseteq [n] \\\\ J_1 \cup \cdots \cup J_r = [n]}}
  \Delta(f_1; x; u_{J_1}) \cdots \Delta(f_r; x; u_{J_r}).
\]

\begin{enumerate}
\def\labelenumi{\arabic{enumi})}
\setcounter{enumi}{1}
\tightlist
\item
  \emph{Boolean Taylor product rule:}
\end{enumerate}

\[
  (f_1 \cdots f_r)(x + {\textstyle\sum_{i=1}^n} u_i)
  =
  \sum_{J_1, \dots, J_r \subseteq [n]}
  \Delta(f_1; x; u_{J_1}) \cdots \Delta(f_r; x; u_{J_r}).
\]

The sum in 1) runs over all \emph{ordered coverings} \((J_1,\dots,J_r)\)
of \([n]\): the \(J_a\) may overlap, and empty \(J_a\) are allowed. Part
2) is the product of Taylor expansions. The binomial versions follow by
applying 1,2) to \(S(\gamma)\).

\end{theorem}

\begin{proof}

By discrete Taylor \ref{prop:taylor-duality},
\(f_a(x + \sum_{i \in S} u_i) = \sum_{J_a \subseteq S} \Delta(f_a; x; u_{J_a})\).
Substituting into the alternating sum for \(f_1 \cdots f_r\) and
expanding the product gives:

\[
  \Delta(f_1 \cdots f_r; x; u_1, \dots, u_n)
  =
  \sum_{\substack{J_1, \dots, J_r \subseteq [n] \\\\ S \supseteq J_1 \cup \cdots \cup J_r}}
  (-1)^{n-|S|}\,
  \Delta(f_1; x; u_{J_1}) \cdots \Delta(f_r; x; u_{J_r}).
\]

Exchanging the order of summation, a fixed tuple
\((J_1, \dots, J_r) \in \KP([n])^r\) appears in the \(S\)-summand
exactly when \(J_1 \cup \cdots \cup J_r \subseteq S\). Setting
\(M = [n] \setminus (J_1 \cup \cdots \cup J_r)\), the sign sum is
\(\sum_{R \subseteq M} (-1)^{|M|-|R|} = (1-1)^{|M|}\): this is \(1\) if
\(J_1 \cup \cdots \cup J_r = [n]\) and \(0\) otherwise.

\end{proof}

\end{document}
